\section{Physical one-step complexity and homogeneous distortion}
\label{sec:uc-complexity}
We verify the remaining curve-growth hypotheses for precisely the class
constructed above. The argument concerns actual one-flight branches and
all their homogeneous pieces. It is not an estimate obtained by counting
long collision words. Constants may depend on the fixed original geometric
family and the chosen finite curvature allowance $B_{\rm cv}$, but not on $R$,
incidences, clearance marks, collision separation or exact-event selectors.

\subsection{Uniform finite geometry and a one-sided first-hit lemma}
The retained $\tau_+$ can be defined intrinsically as the supremum of
lengths of straight free segments in the original family. For clarity, its
finiteness also follows directly from the fixed lattice. A complete line
of irrational direction is dense on $\R^2/\Lambda$ and meets the open
radius-$R_-$ disk. A line of rational direction is parallel to a primitive
lattice vector $z$, with neighboring lattice rows at distance
$(\sqrt3/2)/|z|\le\sqrt3/2$. Its distance from one such row is at most
$\sqrt3/4<R_-$. It therefore meets open disks on that row. The same
argument applies to half-lines. The direction dichotomy follows, for
example, by describing the closure of a line subgroup of a two-dimensional
torus: a proper connected subgroup has an integral annihilator and is a
rational circle; otherwise the closure is the full torus. The positive
semiorbit has the same closure by recurrence of translation on this compact
group. These elementary observations use the Euclidean triangular lattice,
not a linear deformation of the physical disks.

If free segments of unbounded length existed, translate their midpoints
into a compact fundamental cell and take a convergent subsequence of
midpoints and directions. Every finite subsegment of the limiting complete
line would avoid every open radius-$R_-$ disk, a contradiction. Thus
$\tau_+<\infty$, uniformly also for $R\ge R_-$. This compactness proof
does not assign a sampled numerical value to $\tau_+$.
Put
\begin{equation}\label{eq:uc-candidates}
 D_f=\tau_++2R_+,\qquad
 N_f=\#\{z\in\Lambda:0<|z|\le D_f\},\qquad J_f=2N_f+1.
\end{equation}
The packing bound $N_f+1\le4(D_f+1/2)^2$ is fully geometric.

Choose a preliminary curve length small enough that the relative angle
$\eta_z=\arg(z-q)-\vartheta$ varies by at most $\pi/2$ for every
candidate. This is possible uniformly because
$|z-q|\ge1-R_+$ and $\vartheta_s=\kappa+V\le k_++U$.
On every such graph,
\[
 \eta_z'\le(1-R_+)^{-1}-2k_-<0,
 \qquad b_z'=-c-t_z(\kappa+V),
 \quad t_z=(z-q)\cdot v,\quad b_z=(z-q)\cdot v^\perp.
\]
The positive-foot set is an interval, and $b_z$ is strictly decreasing
there. Every disk-intersection set is consequently an interval, with at
most two tangency endpoints. The initial point is exterior to every other
disk, so $t_z>0$ and $|b_z|<R$ are exactly the conditions for a positive
intersection. Distinct first-entry times cannot exchange their order
without a common point of two disjoint disks. Cutting at these at most
$2N_f$ tangencies gives at most $J_f$ genuine first-hit branches.
Coordinate seams are glued on the collision circle, not counted as new
physical branches. In a seam-free local lift the image is a graph and
$\varphi_1$ is strictly monotone by Lemma~\ref{lem:uc-transform}.
Each first-hit branch thus meets each target homogeneity strip at most once.

\begin{lemma}[At most one non-grazing limiting branch]
\label{lem:uc-central}
Fix $k_0\ge2$, with the strips of Theorem~\ref{thm:ag-family}.
There is a common $\ell_0(k_0)>0$ such that every allowed curve of length
at most $\ell_0(k_0)$ has at most one connected physical image piece in
$H_0$. This assertion counts maximal physical homogeneous pieces before
the length subdivision used in the growth partition.
\end{lemma}
\begin{proof}
We give the local argument also at a limiting initial grazing state.
Fix a limiting radius and outgoing state in the compact closure of the
collision cylinder. Consider its oriented straight ray and the first disk
which it intersects strictly. Such a disk exists by the preceding finite
geometry, even if the ray is tangent to one or more earlier disks.
The strict first-entry time and center persist continuously nearby.
The finitely many earlier possible contacts in the limit are tangencies.
Any nearby first-hit branch on one of these earlier disks has terminal
incidence tending to zero, hence cannot have its image in $H_0$ in a
sufficiently small neighborhood. A later disk cannot be the next collision
when the persistent strict disk is met.

It remains to check connectivity for the persistent strict disk, rather
than infer it just from uniqueness of a limiting center. At an earlier
tangency $b_z=\pm R$ the positive foot satisfies
$t_z\ge\sqrt{(1-R_+)^2-R_+^2}>0$. Thus $b_z'$ is uniformly negative
there on all allowed graphs. In a sufficiently small neighborhood this
candidate can cross only one of its two tangency levels: the levels are
separated by $2R\ge2R_-$ and $|b_z'|$ has a uniform finite upper bound
on the finite candidate set. Its intersection condition is therefore
one-sided along any short source graph. Avoiding all earlier disks is an
intersection of one-sided intervals, hence an interval. The strict-disk
branch on that interval has monotone $\varphi_1$, so restricting it to
$H_0$ is still an interval. Incipient tangencies after the strict contact
have no effect on which collision is first.

These neighborhoods are open in radius and the compact closed collision
cylinder. A finite subcover, and a smaller common curve length, give
$\ell_0(k_0)$; the bounded cone slopes keep each sufficiently short graph
inside a member of this cover. Equivalently a sequence of shrinking
counterexample curves would converge to one of the limiting states just
analyzed. Only this limiting line is allowed to pass through tangency
points for purposes of classification. Every branch being estimated is
an actual regular billiard flight with its actual first collision.
No tangency is erased from a physical record.
\end{proof}

\subsection{The complete one-step weighted sum}
For a tangent vector to a positive or negative cone graph define
\[
 |v|_* =\frac{\kappa+|V|}{\sqrt{1+V^2}}|v|.
\]
On these cones $|v|\le|v|_*\le\Gamma|v|$, where
$\Gamma=\sqrt{1+k_+^2}$. This follows respectively from $\kappa>1$
and Cauchy--Schwarz. Extend this cone norm to an equivalent norm off the
cones if needed; no subsequent estimate uses such an extension.
For a positive cone vector the exact expansion is
\begin{equation}\label{eq:uc-adapted}
 \frac{|DT_Rv|_*}{|v|_*}
   =1+\frac{2\kappa A}{dw}
   \ge1+\frac{2k_-t_*}{d}\ge\Lambda_*:=1+2k_-t_*=\frac{59}{47}.
\end{equation}
Time reversal gives the corresponding backward expansion on negative cones.

\begin{theorem}[Uniform one-step expansion and weakened exponent]
\label{thm:uc-expansion}
Choose $k_0$ so large that
\begin{equation}\label{eq:uc-k0}
 T_0:=\frac{J_f}{k_-t_*(k_0-1)}<\frac{1-\Lambda_*^{-1}}2.
\end{equation}
For $\ell_*>0$ sufficiently small as in Lemma~\ref{lem:uc-central},
let $\mathcal H(C)$ be the maximal regular homogeneous pieces of $T_RC$.
Then
\begin{equation}\label{eq:uc-onestep}
 \sum_{V\in\mathcal H(C)}\sup_V J^*_V(T_R^{-1})
       \le\theta_*:=\Lambda_*^{-1}+T_0<1.
\end{equation}
For every $\zeta>1/2$ there is a finite uniform $C_\zeta$ such that
\begin{equation}\label{eq:uc-weakened}
 \sum_{V\in\mathcal H(C)}\bigl(\sup_V J_V(T_R^{-1})\bigr)^\zeta
                         \le C_\zeta.
\end{equation}
Here $J^*$ and $J$ denote curve Jacobians in the adapted and Euclidean
metrics respectively; these are supremum bounds, not just average bounds.
\end{theorem}
\begin{proof}
The possible $H_0$ piece contributes at most $\Lambda_*^{-1}$.
On $H_k^\pm$, $d\le k^{-2}$, so the inverse of
\eqref{eq:uc-adapted} is at most $(2k_-t_* k^2)^{-1}$.
There are at most $J_f$ pieces in each such strip. Sum both signs and use
$\sum_{k\ge k_0}k^{-2}\le1/(k_0-1)$ to prove
\eqref{eq:uc-onestep}. This sum includes every first-hit branch.
The norm comparison gives $J\le\Gamma J^*$ at the two endpoints,
not a new factor for each subsequent collision. Thus one may take
\[
 C_\zeta=(\Gamma/\Lambda_*)^\zeta+
 2J_f\left(\frac{\Gamma}{2k_-t_*}\right)^\zeta
                         \sum_{k\ge k_0}k^{-2\zeta}<\infty.
\]
These are precisely the needed expansion/complexity properties, with
weakened-exponent threshold $\zeta_0=1/2<1$. No claim of the smaller
infinite-horizon threshold appearing elsewhere in the literature is used.
\end{proof}

\subsection{A direct distortion estimate in the target coordinate}
\begin{lemma}[One-step logarithmic distortion]\label{lem:uc-logjac}
Let $J_C T_R$ be the Euclidean curve Jacobian. On a smooth image graph,
\begin{align}
 \partial_{s_1}\log J_C T_R
 &=\frac{p_1w}{A}+\frac{dp(\kappa+2V)}{A^2}
       -\frac{d\tau V_s}{A^2}+\frac{p_1V_1}{d}\notag\\
 &\hspace{8mm}+\frac{V_1 V_{1,s_1}}{1+V_1^2}
       +\frac{VV_s d}{A(1+V^2)}.\label{eq:uc-logjac}
\end{align}
There is a finite uniform $D_B$ such that its magnitude is at most
$U/d+D_B$. For example, with $a_0=2k_-t_*$ and
$B_2=B_{\rm cv}(1+U^2)^{3/2}$, one may use
\begin{equation}\label{eq:uc-DB}
 D_B=t_*^{-1}+2/(t_*a_0)
                 +(B_2/a_0)(1+(2k_-)^{-1})+B_2.
\end{equation}
\end{lemma}
\begin{proof}
The Jacobian is
$J_C T_R=(A/d)\sqrt{1+V_1^2}/\sqrt{1+V^2}$.
Differentiate its logarithm in $s$ and divide by $-A/d$.
Use the exact $A_s,d_s$ of Lemma~\ref{lem:uc-transform}; their apparent
$1/d$ singularity leaves exactly $p_1V_1/d$ and the finite terms in
\eqref{eq:uc-logjac}. This proves the identity, including its signs.
For the bound use $A\ge t_*w\ge a_0$, $w/A\le1/t_*$,
$\kappa+2V\le2w$, $\tau/A\le1/(2k_-)$ and
$|V_s|,|V_{1,s_1}|\le B_2$. The six terms are respectively bounded
by $t_*^{-1}$, $2/(t_*a_0)$, $B_2/(2k_-a_0)$, $U/d$,
$B_2$ and $B_2/a_0$. There is no inverse initial incidence in this
bound. Since $ds_1/d\ell_1\le1$ in absolute value, it also bounds
the logarithmic derivative per unit image arclength.
\end{proof}

\begin{theorem}[Uniform homogeneous distortion at every collision count]
\label{thm:uc-distortion}
There is $D<\infty$ such that, if $x,x'$ belong to a single regular
curve segment whose images through time $l$ are homogeneous, then
\begin{equation}\label{eq:uc-distortion}
 |\log J_C T_R^l(x)-\log J_C T_R^l(x')|
 \le D\,d_{T_R^l C}(T_R^lx,T_R^lx')^{1/3}.
\end{equation}
Distances are along the same physical curve branch. The constant is
uniform in $l$ and $R$. The corresponding Jacobian-ratio distortion bound
and the stable inverse version hold on short homogeneous curves.
\end{theorem}
\begin{proof}
For two points of an image graph in $H_k^\pm$, let $r$ be their
arclength separation. Since $V_1\ge k_-$,
\[
 r\le \frac{\sqrt{1+U^2}}{k_-}
           (k^{-2}-(k+1)^{-2})\le L_h k^{-3},
 \qquad L_h=3\sqrt{1+U^2}/k_-.
\]
Also $d\ge (2/\pi)(k+1)^{-2}$ throughout that segment. The
logarithmic derivative bound is therefore at most $C k^2$, with a
constant independent of $k$. Hence its integral is bounded by
$C k^2r\le C L_h^{2/3}r^{1/3}$.
On $H_0$, $d\ge\sin(k_0^{-2})$ and the same conclusion follows
from the Lipschitz bound, after fixing the maximal short-curve length.
This proves one-step logarithmic distortion in the image distance,
with a common constant $D_1$.

Let $r_l=d_{T_R^l C}(T_R^lx,T_R^lx')$.
The adapted expansion and its endpoint metric comparison imply
$r_j\le\Gamma\Lambda_*^{-(l-j)}r_l$. Expand the logarithm of the
$l$-step Jacobian as the sum of its actual one-step logarithms and apply
the preceding estimate at image times $j+1$. The geometric series gives
\[
 D=D_1\Gamma^{1/3}\sum_{r=0}^{\infty}\Lambda_*^{-r/3}
        =\frac{D_1\Gamma^{1/3}}{1-\Lambda_*^{-1/3}}.
\]
For bounded $r_l$, the inequality $|e^z-1|\le e^{|z|}|z|$ gives
the ratio form. Time reversal $\mathcal I(s,\varphi)=(s,-\varphi)$
is an arclength isometry and conjugates $T_R$ to $T_R^{-1}$, yielding
the stable version. The smooth collision measure is
$c\,ds\,d\varphi/(4\pi R)$ and is invariant, so its measure Jacobian
is identically one. The measure-Jacobian portion of the distortion
hypothesis therefore has zero distortion. This use of time reversal
changes neither the original section nor its occupation convention.
\end{proof}

